\documentclass[12pt]{article}

\usepackage[dvips,letterpaper,margin=0.75in,bottom=0.5in]{geometry}
\usepackage{cite}
\usepackage{slashed}
\usepackage{graphicx}
\usepackage{amsmath}
\usepackage{amssymb}
\usepackage{braket}
\usepackage{pythonhighlight}
\usepackage{mathtools}
\begin{document}

\newcommand{\ihbar}{\ensuremath{i \hbar}}
\newcommand{\Pss}{\ensuremath{\Psi^*}}
\newcommand{\dPsidt}{\ensuremath{ \frac{\partial \Psi}{\partial t} }}
\newcommand{\dPsidx}{\ensuremath{ \frac{\partial \Psi}{\partial x} }}
\newcommand{\ddPsidx}{\ensuremath{ \frac{\partial^2 \Psi}{\partial x^2} }}
\newcommand{\dPssdt}{\ensuremath{ \frac{\partial \Psi^*}{\partial t} }}
\newcommand{\dPssdx}{\ensuremath{ \frac{\partial \Psi^*}{\partial x} }}
\newcommand{\ddPssdx}{\ensuremath{ \frac{\partial^2 \Psi^*}{\partial x^2} }}

\newcommand{\dphidt}{\ensuremath{ \frac{d \phi}{dt} }}
\newcommand{\dpsidx}{\ensuremath{ \frac{d \psi}{dx} }}
\newcommand{\ddpsidx}{\ensuremath{ \frac{d^2 \psi}{dx^2} }}


\title{PHY 115L \\ Final Project Guidelines }
\author{Michael Mulhearn}

\maketitle

\section{Overview}

The final project is an opportunity for your to explore a physics
problem of your own choosing and apply the techniques of computational
physics to solve, illustrate, or in some other way illuminate it.  It
can be connected to techniques or topics we studied in class, or it
can strike out in a completely new direction.  Your problem should be
connected to quantum mechanics or classical mechanics, but I will
interpret this scope quite broadly.  The most important thing is that
your topic is well choosen as something of interest to you, on which
you can make reasonable progress.

In terms of size, this should be something like $50\%$ of the scope of a typical homework.  (This is because the homework lay out a very clear path, so progress is much more rapid.)

\section{Examples}

You are not required to take one of these examples, but here are some examples of reasonably scoped projects:
\begin{itemize}
\item Adding spin-orbit coupling to our numerical model for the neutron.
\item Animating the time evolution of a wave packet in a quantum harmonic oscillator.
\item Animating scattering of a wave packet.
\item Plotting the 3-D solutions to the hydrogen atom in an illustrative manner.
\item Exploring the ``shooting method'' to finding energy eigenstates.
\item Develop a tool for calculating (exactly) the Clebsch-Gordan coefficients.
\end{itemize}

\end{document}
